\section{Thiết lập và vận hành nhóm}
% [NHÁP -- Chủ trì: Nguyễn Văn Thương (quản lý) & Vương Đình Minh Trí (DevOps).]

Nhóm vận hành trên bốn công cụ, mỗi công cụ một vai trò rõ ràng, kết nối bằng đồng bộ tự động để tránh nhập tay trùng lặp:
\begin{itemize}
    \item \textbf{JIRA} (key SCRUM) -- điểm tạo và quản lý công việc (Scrum board, sprint): \url{https://thng.atlassian.net/jira/software/projects/SCRUM/boards/1}.
    \item \textbf{Notion} -- Single Source of Truth: tài liệu, kế hoạch, ADR, log thí nghiệm, knowledge hub (không gian làm việc nội bộ của nhóm).
    \item \textbf{Discord} (server ID \texttt{1511918814995157074}) -- kênh trao đổi và nhận thông báo tự động.
    \item \textbf{GitHub} -- mã nguồn, Pull Request, review (\texttt{Disease-Diagnosis-RAG-System}).
\end{itemize}

\begin{figure}[H]
    \centering
    \resizebox{0.95\textwidth}{!}{%
        \begin{tikzpicture}
            % --- Nút đặt bằng toạ độ tuyệt đối để tránh chồng lấp ---
            \node[stage, text width=2.6cm] (notion)  at (0,1.5)   {\textbf{Notion}\\Source of Truth};
            \node[stage, text width=2.6cm] (jira)    at (-5.2,1.5) {\textbf{JIRA}\\Scrum board};
            \node[stage, text width=2.6cm] (discord) at (5.2,1.5)  {\textbf{Discord}\\Thông báo nhóm};
            \node[stage, text width=2.6cm] (github)  at (0,3.7)   {\textbf{GitHub}\\Code + PR};
            \node[future, text width=2.6cm] (prism)   at (0,-0.5)  {\textbf{Prism/Overleaf}\\Báo cáo};
            % --- Mũi tên ---
            \draw[twoway] (jira) -- node[above, font=\footnotesize, align=center]{Unito\\2 chiều} (notion);
            \draw[flow] (github) -- node[right, font=\footnotesize]{evidence/PR} (notion);
            \draw[flow] (notion) -- node[right, font=\footnotesize]{nội dung} (prism);
            % Kênh thông báo JIRA -> Discord đi vòng dưới (Manhattan), nằm hẳn dưới Prism
            \draw[flow] (jira.south) |- (0,-2.6) -| (discord.south);
            \node[font=\footnotesize, align=center] at (0,-2.95) {Jira Automation, 1 chiều};
        \end{tikzpicture}}
    \caption{Sơ đồ đồng bộ công cụ: JIRA $\leftrightarrow$ Notion (2 chiều, Unito), JIRA $\rightarrow$ Discord (1 chiều, Jira Automation). Notion và Discord không đồng bộ trực tiếp.}
    \label{fig:toolchain}
\end{figure}

\subsection{Đồng bộ JIRA $\leftrightarrow$ Notion (Unito, 2 chiều)}
\begin{table}[H]
    \centering
    \small
    \caption{Sáu trường đồng bộ hai chiều giữa JIRA và Notion (Unito).}
    \label{tab:unito}
    \begin{tabular}{lll}
        \hline
        \textbf{JIRA}        & \textbf{Notion} & \textbf{Chiều}     \\
        \hline
        Summary              & Task Name       & 2 chiều            \\
        Status               & Status          & 2 chiều            \\
        Priority             & Priority        & 2 chiều            \\
        Sprint               & Sprint          & 2 chiều            \\
        Due date             & Deadline        & 2 chiều            \\
        Description (footer) & Link            & 1 chiều (gắn link) \\
        \hline
    \end{tabular}
\end{table}

\subsection{Thông báo JIRA $\rightarrow$ Discord (Jira Automation, 1 chiều)}
Sáu luồng thông báo tự động gửi qua webhook vào kênh Discord (Webhook \#1 và \#2 theo loại sự kiện):
\begin{enumerate}
    \item Tạo issue mới (issue created).
    \item Đổi trạng thái (status changed).
    \item Đổi người phụ trách (assignee changed).
    \item Thêm bình luận (comment added).
    \item Bắt đầu / kết thúc sprint.
    \item Hoàn thành issue (done).
\end{enumerate}
\textit{Lưu ý bảo mật:} không in URL webhook đầy đủ trong tài liệu; chỉ cấu hình trong JIRA Automation.

\textbf{Nguyên tắc:} Notion là nguồn sự thật cho tài liệu; JIRA là nguồn sự thật cho công việc; mọi task chỉ tạo ở một nơi (JIRA) rồi đồng bộ sang Notion; Discord chỉ nhận thông báo một chiều. Phiên bản JIRA miễn phí giới hạn 10 người dùng; quản trị: Nguyễn Văn Thương.
